\subsection{隐藏通道熵影、双宿主边界 faithful carrier、有限 jet 刚性与 Schur--median 终端律}\label{subsec:conclusion-hiddenchannel-boundarycarrier-finitejet-schur-median-rigidity}

沿用定理 \ref{thm:conclusion-chebotarev-hiddenchannel-dominant-entropy-barrier}、定理 \ref{thm:conclusion-hidden-artin-ghost-oscillation}、命题 \ref{prop:xi-time-part9g-first-deviation-depth}、定理 \ref{thm:xi-time-part9g-langlands-euler-separation}、定理 \ref{thm:conclusion-boundary-godel-adhesion-defect-phase-diagram}、定理 \ref{thm:conclusion-boundary-godel-bitlength-exponent-equals-adhesion}、定理 \ref{thm:conclusion-prime-log-abel-mellin-pole-rigidity}、推论 \ref{cor:conclusion-prime-log-alpha-singular-rank-one}、定理 \ref{thm:conclusion-prime-log-rh-operator-countable-prime-rank}、定理 \ref{thm:conclusion-hankel-arithmetic-tomography-finite-bad-primes-completion}、定理 \ref{thm:conclusion-hankel-stiffness-dual-role}、定理 \ref{thm:derived-chain-arithmetic-median-unique-minimizer} 与推论 \ref{cor:derived-chain-arithmetic-boolean-interval-geodesics} 的记号。前述结果已经分别给出：primitive Chebotarev 群商熵的隐藏主导障碍与幽灵统计、停机型二进扭子偏离的首次显影深度、subfull regime 中边界复杂度对粘连缺陷的完全塌缩、prime-log 校准的单极点边界方向与可数无限 prime-register 宿主、有限 Hankel realization 的好素数补全与 \(p\)-进刚性，以及有限链内算子的算术中值与布尔测地计数。把这些接口并读后，还可继续压缩为如下七条终端后果。

\leanverified{paper\_conclusion\_hiddenchannel\_square\_root\_square\_law}
\begin{theorem}[隐藏主导字符的平方根--平方律]\label{thm:conclusion-hiddenchannel-square-root-square-law}
设
\[
\pi:G_2\twoheadrightarrow G_1,\qquad K:=\ker \pi,
\]
并固定 \(u\in(0,1]\)。假设隐藏谱中存在唯一主导字符
\[
\chi_\ast\in \widehat G_2\setminus K^\perp
\]
及其简单特征值 \(\mu_\ast\)，满足
\[
|\mu_\ast|=\eta_{\mathrm{hid}}(u,\pi),
\]
且该通道在 \(\mu_\ast\) 处具有谱隙。则存在零均值类函数
\[
\varepsilon:G_2\to\CC,
\qquad
\pi_\ast\varepsilon=0,
\]
以及一条无穷子序列 \(n_j\to\infty\)，使得
\[
\Delta_{n_j}
\asymp
\frac{n_j^2}{\lambda_1(u)^{2n_j}}
\bigl|\Pi_{n_j}(\varepsilon;u)\bigr|^2,
\]
其中
\[
\Pi_n(\varepsilon;u)
:=
\sum_{c_2\in G_2}\varepsilon(c_2)\,B_{n,c_2}^{(2)}(u),
\]
\[
\Delta_n
:=
D_{\mathrm{KL}}\!\bigl(\mu_{n,u}^{(2)}\Vert u_{G_2}\bigr)
-
D_{\mathrm{KL}}\!\bigl(\mu_{n,u}^{(1)}\Vert u_{G_1}\bigr).
\]
等价地，在线性统计层，最慢隐藏通道以
\[
\Pi_n(\varepsilon;u)\sim c_\varepsilon \mu_\ast^n/n
\]
的幽灵振荡存活；而在熵层，同一通道只以平方速率留下阴影：
\[
\Delta_n\asymp
\left(\frac{\eta_{\mathrm{hid}}(u,\pi)}{\lambda_1(u)}\right)^{2n}.
\]
故群商熵损失与隐藏幽灵线性统计之间存在严格的平方根--平方对应。
\end{theorem}

\begin{proof}
由定理 \ref{thm:conclusion-hidden-artin-ghost-oscillation}，存在零均值类函数 \(\varepsilon\) 与常数 \(c_\varepsilon\neq 0\) 使
\[
\Pi_n(\varepsilon;u)
=
c_\varepsilon\frac{\mu_\ast^n}{n}
+
O\!\left(\frac{\eta_-^n}{n}\right),
\qquad
\eta_-<\eta_{\mathrm{hid}}(u,\pi).
\]
于是沿某条无穷子序列 \(n_j\to\infty\) 有
\[
\bigl|\Pi_{n_j}(\varepsilon;u)\bigr|^2
=
|c_\varepsilon|^2\eta_{\mathrm{hid}}(u,\pi)^{2n_j}n_j^{-2}(1+o(1)).
\]
另一方面，由定理 \ref{thm:conclusion-chebotarev-hiddenchannel-dominant-entropy-barrier}，
\[
\Delta_{n_j}
\asymp
\left(\frac{\eta_{\mathrm{hid}}(u,\pi)}{\lambda_1(u)}\right)^{2n_j}.
\]
将两式比较即得
\[
\Delta_{n_j}
\asymp
\frac{n_j^2}{\lambda_1(u)^{2n_j}}
\bigl|\Pi_{n_j}(\varepsilon;u)\bigr|^2.
\]
证毕。
\end{proof}

\leanverified{paper\_conclusion\_halting\_walsh\_plateau\_localfactor\_locking}
\begin{theorem}[Walsh 首断点、账本平台与局部因子误差壁垒的刚性锁定]\label{thm:conclusion-halting-walsh-plateau-localfactor-locking}
设 \(B=2^L\)。对任意图灵机--输入对 \((M,x)\)，定义停机编码
\[
\alpha_{M,x}:=\sum_{n\ge 0} t_n(M,x)B^n\in \ZZ_2,
\]
并令
\[
N\in\NN\cup\{\infty\}
\]
为首次停机时间。记 \(k^\ast(M,x)\) 为相应 Walsh 审计中的首个非平凡断点，\((\kappa_m)_{m\ge 0}\) 为账本缺陷剖面，并定义平台起点
\[
m_{\mathrm{plat}}(M,x)
:=
\inf\{\,m\ge 0:\ \kappa_{m'}=\kappa_m\ \text{对全部 }m'\ge m\,\}.
\]
则若 \(N<\infty\)，有精确恒等式
\[
m_{\mathrm{plat}}(M,x)=N=\frac{k^\ast(M,x)-1}{L}.
\]
并且对任意固定 \(s>0\) 与素数 \(p\)，局部因子误差满足
\[
\bigl|L_p(s,g_{M,x})-(1-p^{-s})^{-2}\bigr|
\le
C(s,p)\,2^{-N}
=
C(s,p)\,2^{-(k^\ast(M,x)-1)/L}.
\]
若 \(N=\infty\)，则
\[
k^\ast(M,x)=\infty,\qquad
m_{\mathrm{plat}}(M,x)=\infty,\qquad
g_{M,x}=I_2,
\]
从而局部因子严格退化为平凡模型。
因此，首次停机时间同时等于二进赋值深度、Walsh 首断点、账本平台起点与局部 Euler 因子接近平凡模型的指数尺度。
\end{theorem}

\begin{proof}
由停机型 \(2^L\)-进编码链已经证明的等价关系，
\[
N<\infty
\iff
\alpha_{M,x}=-B^N
\iff
v_2(\alpha_{M,x})=LN
\iff
k^\ast(M,x)=LN+1
\iff
\kappa_m=\min(m,N)\log 2.
\]
故当 \(N<\infty\) 时，\((\kappa_m)\) 的平台恰从 \(m=N\) 开始，即
\[
m_{\mathrm{plat}}(M,x)=N=\frac{k^\ast(M,x)-1}{L}.
\]
非停机时 \(\alpha_{M,x}=0\)，于是上述四个量都退化到 \(\infty\) 或平凡模型。

另一方面，命题 \ref{prop:xi-time-part9g-first-deviation-depth} 把首次偏离深度精确识别为 \(LN+1\)，而定理 \ref{thm:xi-time-part9g-langlands-euler-separation} 给出停机深度为 \(N\) 时的 Euler 型分离率
\[
\bigl|L_p(s,g_{M,x})-(1-p^{-s})^{-2}\bigr|
\le
C(s,p)\,4^{-N/2}
=
C(s,p)\,2^{-N}.
\]
再代入
\(
N=(k^\ast(M,x)-1)/L
\)
即得所示公式。证毕。
\end{proof}

\begin{theorem}[边界 faithful carrier 的双宿主分裂]\label{thm:conclusion-boundary-faithful-carriers-double-host-splitting}
不存在任何有限秩 additive/localized ledger \(H\)，能够同时 faithful 地承载下列两类数据：
\begin{enumerate}
  \item 全部 subfull 维 regime \((d_G(A)<n)\) 中的边界复杂度；
  \item 全部 prime-log 校准边界流 \(\Psi_\alpha\) \((\alpha>0)\) 及其 RH 候选宿主语义。
\end{enumerate}
更精确地说，在 \(d_G(A)<n\) 时有
\[
d_{\partial G}(A)=\gamma_{\mathrm{adh}}(A),
\]
且精确边界 Gödel 码长指数同样等于 \(\gamma_{\mathrm{adh}}(A)\)；另一方面，prime-log 校准链中的 \(\alpha\)-依赖只以一维边界留数
\[
c\,\zeta(s)
\]
的形式出现，而 faithful 宿主必须满足
\[
\ell_{\mathrm{loc}}(M_\alpha)=\aleph_0.
\]
于是任何共同 faithful 的最小宿主都必严格分裂为
\[
H_{\mathrm{geom}}\oplus H_{\mathrm{arith}},
\]
其中 \(H_{\mathrm{geom}}\) 只承载边界 Gödel/粘连缺陷数据，而 \(H_{\mathrm{arith}}\) 必包含可数无限 prime-register 轴。故边界 faithful carrier 的最小原则不是单宿主原则，而是双宿主原则。
\end{theorem}

\begin{proof}
在 subfull regime 中，定理 \ref{thm:conclusion-boundary-godel-adhesion-defect-phase-diagram} 给出
\[
d_{\partial G}(A)=\gamma_{\mathrm{adh}}(A),
\]
而定理 \ref{thm:conclusion-boundary-godel-bitlength-exponent-equals-adhesion} 进一步给出精确边界 Gödel 码位长的指数同样等于 \(\gamma_{\mathrm{adh}}(A)\)。因此，此时几何边界复杂度完全塌缩到粘连缺陷单轴上。

另一方面，定理 \ref{thm:conclusion-prime-log-abel-mellin-pole-rigidity} 给出
\[
\mathcal D_\alpha(s)=c\,\zeta(s)+\mathcal J_\alpha(s),
\]
其中 \(\mathcal J_\alpha\) 在 \(\Re s>0\) 上全纯；推论 \ref{cor:conclusion-prime-log-alpha-singular-rank-one} 进一步说明 \(\alpha\) 的全部非全纯依赖都被压缩到 \(\CC\cdot \zeta(s)\) 这一维边界方向。最后，由定理 \ref{thm:conclusion-prime-log-rh-operator-countable-prime-rank}，任意 faithful 承载 prime-log 校准的 RH 候选宿主必须满足
\[
\ell_{\mathrm{loc}}(M_\alpha)=\aleph_0.
\]
故前一类数据要求的是有限几何粘连宿主，而后一类数据要求的是可数无限 prime-register 宿主，二者不可能在同一个有限秩账本中同时 faithful 地闭合。于是最小 faithful carrier 必分裂为几何/算术双宿主。证毕。
\end{proof}

\begin{theorem}[有限 jet 的 \(p\)-进刚性与 Archimedean 余维一塌缩]\label{thm:conclusion-finitejet-goodprime-archimedean-codimone-collapse}
在稳定 defect-\(\kappa\) regime 中，固定有限证书
\[
T_\kappa:=(u_0,\dots,u_{2\kappa-1};p),
\]
其中 \(p\) 为一枚好素数，并满足
\[
p\nmid D\det(H_0),\qquad
H_0:=(u_{i+j})_{0\le i,j\le \kappa-1},
\qquad
H_1:=(u_{i+j+1})_{0\le i,j\le \kappa-1}.
\]
记
\[
A:=H_0^{-1}H_1,\qquad
\epsilon:=\kappa_{\mathrm{ind}}-S_{\mathrm{def}}.
\]
则：
\begin{enumerate}
  \item 有限证书 \(T_\kappa\) 唯一决定
  \[
  H_0,\ H_1,\ A,\ \{z_j\}_{j=1}^{\kappa},\ \{a_j\}_{j=1}^{\kappa},\ \det P,\ S_{\mathrm{gen}}.
  \]
  \item 一切算术模糊性都严格局域于有限坏素数集
  \[
  S_{\mathrm{bad}}:=\{p:\ p\mid \det(H_0)\}.
  \]
  对任意 \(p\notin S_{\mathrm{bad}}\) 与任意 \(k\ge 1\)，同余系统
  \[
  H_0X\equiv H_1 \pmod{p^k}
  \]
  有唯一解，且
  \[
  X\equiv A \pmod{p^k}.
  \]
  \item 在 Archimedean 侧，若定义阈值
  \[
  \theta_\epsilon:=4\pi\kappa e^{-1}\epsilon,
  \]
  则
  \[
  \sigma_j(H_{\kappa-1})\le \theta_\epsilon\qquad (2\le j\le \kappa),
  \]
  且
  \[
  \sigma_1(H_{\kappa-1})\ge 4\pi S_{\mathrm{def}}-\theta_\epsilon.
  \]
  特别地，只要
  \[
  \epsilon< \frac{e\,S_{\mathrm{def}}}{2\kappa},
  \]
  就有
  \[
  \#\{j:\sigma_j(H_{\kappa-1})>\theta_\epsilon\}=1.
  \]
\end{enumerate}
故同一有限 jet 证书同时触发两种互补刚性：在 \(p\)-进方向上，模糊性只剩有限坏素数局域化；在实谱方向上，不确定性被压缩到唯一宏观奇异方向，其余 \(\kappa-1\) 个方向全部线性塌缩。
\end{theorem}

\begin{proof}
有限 jet、单个补样与好素数证书的闭包定理已经表明：\(T_\kappa\) 唯一决定 \(H_0,H_1\) 及其伴随 realization 数据 \(A,\{z_j\},\{a_j\},\det P,S_{\mathrm{gen}}\)，这给出第 \(1\) 点。

对第 \(2\) 点，好素数条件意味着 \(H_0\) 在 \(\ZZ_p\) 中可逆。定理 \ref{thm:conclusion-hankel-stiffness-dual-role} 已表明：当 \(p\nmid \det(H_0)\) 时，最小 realization
\[
A=H_0^{-1}H_1
\]
属于 \(M_\kappa(\ZZ_p)\)，且模 \(p^k\) 数据只能有唯一的 lift；定理 \ref{thm:conclusion-hankel-arithmetic-tomography-finite-bad-primes-completion} 又说明，全部算术失真只能发生在有限坏素数集上。故
\[
H_0X\equiv H_1\pmod{p^k}
\]
对每个 \(p\notin S_{\mathrm{bad}}\) 都有唯一解，且该解恰为 \(A\) 的模 \(p^k\) 约化。

第 \(3\) 点则是稳定 defect-\(\kappa\) 扇区中已建立的奇异值估计：
\[
\sigma_1(H_{\kappa-1})\ge 4\pi S_{\mathrm{def}}-4\pi\kappa e^{-1}\epsilon,
\qquad
\sigma_j(H_{\kappa-1})\le 4\pi\kappa e^{-1}\epsilon\quad (j\ge 2).
\]
把右端统一记为 \(\theta_\epsilon\) 即得结论。若再有 \(\epsilon<eS_{\mathrm{def}}/(2\kappa)\)，则
\[
4\pi S_{\mathrm{def}}-\theta_\epsilon>\theta_\epsilon,
\]
故仅第一奇异值超过阈值 \(\theta_\epsilon\)。证毕。
\end{proof}

\begin{theorem}[均匀权重前沿与精确熵饱和的排斥律]\label{thm:conclusion-defect-uniform-frontier-entropy-saturation-exclusion}
在稳定 defect-\(\kappa\) regime 中，记
\[
\epsilon:=\kappa_{\mathrm{ind}}-S_{\mathrm{def}},
\qquad
V_N:=-\log\prod_{j=1}^{\kappa}|\lambda_j^-|,
\]
并把缺陷权重写为
\[
\{w_\ell\}_{\ell=1}^{\kappa_{\mathrm{ind}}}\subset(0,1),
\qquad
\bar w:=\frac{S_{\mathrm{def}}}{\kappa_{\mathrm{ind}}},
\qquad
\mathrm{Var}(w):=
\frac{1}{\kappa_{\mathrm{ind}}}
\sum_{\ell=1}^{\kappa_{\mathrm{ind}}}(w_\ell-\bar w)^2.
\]
则同时成立：
\begin{enumerate}
  \item 第一谐波满足 sharp 上界
  \[
  |u_1|
  \le
  \frac{e}{4\pi}
  \frac{S_{\mathrm{def}}(\kappa_{\mathrm{ind}}-S_{\mathrm{def}})}{\kappa_{\mathrm{ind}}}
  -\kappa_{\mathrm{ind}}\mathrm{Var}(w).
  \]
  因而在固定 \((\kappa_{\mathrm{ind}},S_{\mathrm{def}})\) 下，\(|u_1|\) 的最大允许值只能由均匀权重
  \[
  \mathrm{Var}(w)=0
  \]
  实现。
  \item 负谱体积势满足 blow-up 下界
  \[
  V_N\ge (2\kappa-2)\log \epsilon^{-1}-O_{\kappa,S_{\mathrm{def}}}(1)
  \qquad(\epsilon\downarrow 0).
  \]
  \item 若 \(\kappa\ge 2\)，则精确熵饱和
  \[
  \kappa_{\mathrm{ind}}=S_{\mathrm{def}}
  \]
  不可能发生；事实上这会强制
  \[
  u_n=0\quad(n\ge 1),\qquad
  H_{\kappa-1}=u_0e_0e_0^\top,
  \]
  从而
  \[
  \rank(H_{\kappa-1})=1,
  \]
  与稳定 defect-\(\kappa\) 的秩条件矛盾。
\end{enumerate}
因此，稳定多缺陷 regime 的可行区域被两条互补前沿夹住：其上边界是均匀权重极值前沿，其下边界是 \((2\kappa-2)\)-阶体积 blow-up 壁，而精确熵饱和本身被严格排除。
\end{theorem}

\begin{proof}
第一条是稳定 defect-\(\kappa\) 口径下的第一谐波方差惩罚上界。其右端在固定 \((\kappa_{\mathrm{ind}},S_{\mathrm{def}})\) 时对 \(\mathrm{Var}(w)\) 严格单调递减，故只有 \(\mathrm{Var}(w)=0\) 的均匀权重才可能实现最大值。

第二条来自负谱乘积上界
\[
\prod_{j=1}^{\kappa}|\lambda_j^-|
\le
C_{\kappa,S_{\mathrm{def}}}\,\epsilon^{2\kappa-2},
\]
取负对数即得
\[
V_N\ge (2\kappa-2)\log \epsilon^{-1}-O_{\kappa,S_{\mathrm{def}}}(1).
\]

对第三条，若 \(\kappa_{\mathrm{ind}}=S_{\mathrm{def}}\)，则已建立的精确熵饱和判据强制全部非零谐波消失，即
\[
u_n=0\quad(n\ge 1),
\]
于是
\[
H_{\kappa-1}=u_0e_0e_0^\top
\]
为秩一矩阵。这与稳定 defect-\(\kappa\) regime 要求
\[
\rank(H_{\kappa-1})=\kappa\ge 2
\]
矛盾。故精确熵饱和不可能发生。证毕。
\end{proof}

\begin{theorem}[统一分离阈值下 Schur 补账本的指数瓶颈]\label{thm:conclusion-schur-ledger-uniform-separation-exponential-bottleneck}
设
\[
A=\{a_1,\dots,a_\kappa\}\subset \DD
\]
为有限缺陷配置，端点 Poisson 权重记为
\[
p_j:=P_{a_j}(-1),
\qquad
p_\Sigma:=\sum_{j=1}^{\kappa}p_j.
\]
令 \((\pi_m)_{m=1}^{\kappa}\) 为相应 Schur 补账本中的步通量，并定义交叉互能
\[
E_{\mathrm{int}}
:=
2\sum_{1\le j<k\le \kappa}\bigl(-\log \rho(a_j,a_k)\bigr).
\]
若存在 \(\rho_0\in(0,1)\) 使
\[
\rho(a_j,a_k)\le \rho_0
\qquad (j\neq k),
\]
则存在某个指标 \(m^\ast\in\{1,\dots,\kappa\}\) 使
\[
\pi_{m^\ast}\le p_\Sigma\,\rho_0^{\,\kappa-1}.
\]
若进一步满足总端点通量的线性控制
\[
p_\Sigma\le C\kappa,
\]
则可 sharpen 为
\[
\pi_{m^\ast}\le C\kappa\,\rho_0^{\,\kappa-1}.
\]
因此，在统一伪双曲分离阈值下，Schur 补账本中必然出现一个指数瓶颈步骤；它的压缩率对缺陷数 \(\kappa\) 呈严格指数衰减，而不是仅仅线性恶化。
\end{theorem}

\begin{proof}
前述 Pick--Poisson/Schur ledger 链已经证明：存在某个 \(m^\ast\) 使
\[
\pi_{m^\ast}\le p_{m^\ast}\exp(-E_{\mathrm{int}}/\kappa).
\]
而在统一分离阈值 \(\rho(a_j,a_k)\le \rho_0\) 下，
\[
E_{\mathrm{int}}
=
2\sum_{1\le j<k\le \kappa}\bigl(-\log \rho(a_j,a_k)\bigr)
\ge
2\binom{\kappa}{2}(-\log \rho_0)
=
\kappa(\kappa-1)(-\log \rho_0).
\]
故
\[
\exp(-E_{\mathrm{int}}/\kappa)\le \rho_0^{\,\kappa-1}.
\]
再用
\[
p_{m^\ast}\le p_\Sigma
\]
便得到
\[
\pi_{m^\ast}\le p_\Sigma\,\rho_0^{\,\kappa-1}.
\]
若再有 \(p_\Sigma\le C\kappa\)，则第二式随之成立。证毕。
\end{proof}

\begin{theorem}[算术中值的唯一极小性与测地退化的阶乘律]\label{thm:conclusion-chain-arithmetic-median-factorial-geodesics}
\leanverified{paper\_conclusion\_chain\_arithmetic\_median\_factorial\_geodesics}
设 \(P\) 为有限链，\(J(P)\) 为其 inner operators 的分配格。取单射素数标号
\[
q:P\hookrightarrow\mathbb P,
\]
并记平方自由算术嵌入为
\[
\Phi_q:J(P)\hookrightarrow \NN_{\mathrm{sf}}.
\]
对任意三点 \(I_1,I_2,I_3\in J(P)\)，记
\[
n_i:=\Phi_q(I_i),
\qquad
I_{123}:=\operatorname{med}_{J(P)}(I_1,I_2,I_3).
\]
则：
\begin{enumerate}
  \item 算术中值满足闭式
  \[
  \Phi_q(I_{123})
  =
  \frac{\gcd(n_1,n_2)\gcd(n_2,n_3)\gcd(n_3,n_1)}
  {\gcd(n_1,n_2,n_3)^2},
  \]
  并且在椭圆族 \((\mathcal E,\delta_{\mathcal E})\) 中，\(E_{\Phi_q(I_{123})}\) 是唯一极小化
  \[
  E_n\longmapsto
  \delta_{\mathcal E}(E_n,E_{n_1})
  +\delta_{\mathcal E}(E_n,E_{n_2})
  +\delta_{\mathcal E}(E_n,E_{n_3})
  \]
  的点。
  \item 若 \(u,v\in \Phi_q(J(P))\) 满足 \(u\mid v\)，并记
  \[
  \operatorname{supp}(v/u)=\{p_1,\dots,p_r\},
  \]
  则 divisibility 区间 \([u,v]_{\mid}\) 与 \(r\)-维布尔立方同构，且从 \(u\) 到 \(v\) 的最短 \(\delta_{\mathcal E}\)-测地条数恰为
  \[
  r!.
  \]
\end{enumerate}
因此，该 faithful 算术模型中，中值唯一性与测地唯一性严格分离：三点中值只有一个，而端点间最短输运却呈现阶乘级退化。
\end{theorem}

\begin{proof}
第 \(1\) 条正是定理 \ref{thm:derived-chain-arithmetic-median-unique-minimizer}：算术中值的闭式由平方自由赋值几何直接给出，而椭圆实现 \(\Psi_q(I)=E_{\Phi_q(I)}\) 保持三元中值与距离，从而把 \(I_{123}\) 送到总距离和的唯一极小点。

第 \(2\) 条则由推论 \ref{cor:derived-chain-arithmetic-boolean-interval-geodesics} 直接给出。对可除区间 \([u,v]_{\mid}\)，支撑差 \(\operatorname{supp}(v/u)=\{p_1,\dots,p_r\}\) 的每个子集都对应一个中间点，故该区间与 \(2^{[r]}\) 同构；而每条最短测地都恰对应于把 \(p_1,\dots,p_r\) 逐个并入的某个排列，故其总数正是 \(r!\)。证毕。
\end{proof}

综上，这一链条把几类原本分散的有限证书刚性焊接到同一终端图景中：hidden Chebotarev 通道把线性幽灵与 KL 熵影锁成平方根--平方律；停机型 \(2^L\)-进偏离把 Walsh 首断点、账本平台与 Euler 局部因子误差锁成同一有限视界深度；boundary faithful carrier 不能压缩为单宿主，而被迫分裂为几何粘连轴与可数 prime-register 轴；而 stable defect-\(\kappa\) 的有限 jet 证书又同时把 \(p\)-进模糊性压到有限坏素数，把 Archimedean 模糊性压到唯一宏观奇异方向，并在统一分离阈值下强制 Schur 补账本出现指数瓶颈。最后，有限链内算子的平方自由算术实现进一步表明：唯一中值与阶乘级测地退化可以在同一 faithful 模型中并存，从而把“唯一恢复”与“唯一输运”彻底区分开来。
